\chapter{Этап 9 — документация Doxygen }
\hypertarget{md__r_e_a_d_m_e}{}\label{md__r_e_a_d_m_e}\index{Этап 9 — документация Doxygen@{Этап 9 — документация Doxygen}}
\label{md__r_e_a_d_m_e_autotoc_md0}%
\Hypertarget{md__r_e_a_d_m_e_autotoc_md0}%


На этом этапе добавлены подробные комментарии Doxygen к шаблонному классу, его специализациям, полям и методам, глобальным шаблонным функциям, функциям загрузки данных и операциям консольного меню.

Комментарии описывают назначение сущностей, шаблонные параметры ({\ttfamily @tparam}), входные параметры ({\ttfamily @param}), возвращаемые значения ({\ttfamily @return}), возможные исключения ({\ttfamily @throws}) и особенности работы ({\ttfamily @note}). Реализация программы не изменялась.

{\bfseries{Сообщение коммита:}}


\begin{DoxyCode}{0}
\DoxyCodeLine{Добавлена\ документация\ Doxygen\ для\ классов\ и\ функций\ проекта}

\end{DoxyCode}


{\bfseries{Комментарий коммита:}}


\begin{DoxyCode}{0}
\DoxyCodeLine{Добавлены\ подробные\ комментарии\ Doxygen\ к\ шаблонному\ классу\ RingBufferCore,\ специализациям\ RingBuffer\ и\ функциям\ проекта.\ Описаны\ параметры\ шаблонов,\ входные\ аргументы,\ возвращаемые\ значения,\ исключения\ и\ особенности\ работы\ методов.\ Документация\ помогает\ разобраться\ в\ назначении\ каждой\ функции\ и\ логике\ кольцевого\ буфера.}

\end{DoxyCode}


clang++ -\/std=c++17 -\/Wall -\/Wextra -\/pedantic -\/Iinclude src/\+main.\+cpp -\/o build/\+\+Ring\+Buffer

Шаблон -\/ функции без чётко заданных типов данных. Типы заданы как параметры. T -\/ параметр типа. Он не известен (не связан ни с каким типом данных).

Специализация -\/ Нетрививально -\/

Во всех пяти функциях происходит практически одинаковое действие:
\begin{DoxyEnumerate}
\item Получаем номер буфера от пользователя.
\item Проверяем, к какой группе относится номер.
\item Переводим общий номер в индекс массива.
\item Вызываем нужную операцию с выбранным объектом
\end{DoxyEnumerate}
\begin{DoxyEnumerate}
\item Допустим, пользователь ввёл номер 7.
\end{DoxyEnumerate}

Программа проверяет диапазоны:
\begin{DoxyItemize}
\item 1–2 — int\+Buffers5;
\item 3–4 — double\+Buffers5;
\item 5–6 — float\+Buffers7;
\item 7–8 — int\+Buffers10;
\item 9–10 — double\+Buffers10. Число 7 попадает в диапазон 7–8. Значит, нужный буфер находится в массиве int\+Buffers10.
\end{DoxyItemize}

Именно это делает условие: else if (index \texorpdfstring{$<$}{<}= int\+Count5 + double\+Count5 + float\+Count7 + int\+Count10)


\begin{DoxyEnumerate}
\item Переводим общий номер в индекс массива Знаем, что буфер находится в int\+Buffers10.
\end{DoxyEnumerate}

Но пользователь ввёл номер 7, а массив начинается с индекса 0. Поэтому нужно преобразовать общий номер в локальный индекс массива.

Перед четвёртой группой находятся:
\begin{DoxyItemize}
\item 2 буфера int\+Buffers5;
\item 2 буфера double\+Buffers5;
\item 2 буфера float\+Buffers7. Всего перед ней 6 буферов.
\end{DoxyItemize}

Поэтому: 7-\/2-\/2-\/2=1

Получается индекс 1, то есть второй элемент массива int\+Buffers10.

В коде: show\+Buffer\+Capacity( int\+Buffers10\+[index -\/ int\+Count5 -\/ double\+Count5 -\/ float\+Count7 -\/ 1] ); 